\scanpage{170}
N.B. Les \unsure{élémts|éléments} de $A$ (formes linéaires continues
sur $U$) sont exactement les formes du type
$u\longrightarrow\langle u.a,b'\rangle$, où $a,b'\in V\in K^{\underline U}$, i.e. des
\unsure{coéfficients|coëfficients|coefficients} de représentation linéaires continues
de dim finie de $\underline U$. \deleted{\unsure{Ces représent.}}

\begin{center}\shortstack{\rule{1.4cm}{0.4pt}\\[-0.9ex]\rule{1.2cm}{0.4pt}}\end{center}

\deleted{Éléments}
\uline{Classes remarquables d'éléments dans $\underline U$}

1) Le groupe $U^*$\deleted{\unsure{$G$}} des éléments inversibles de $U$.
Correspond aux \uline{automorphismes} du foncteur
$I^{\underline U}$, \deleted{qui} i.e. la \deleted{\unsure{donnée pour tout}}\inserted{donnée pour} tout $\unsure{V|W}\in K^{\underline U}$
d'un \deleted{\unsure{automor|automorphisme}} $k$-automorphisme de $V$,
avec condition de commutation pour les
$\underline U$-homomorphismes. \uline{Il y en a beaucoup}
\uline{trop}.

\unsure{3|2}) Le \deleted{\unsure{groupe des élém|groupe des éléments}}\inserted{sous-espace $\unsure{\mathfrak g|g}$ des} éléments primitifs
Il s'agit des endomorphismes du foncteur
additif $I^{\underline U}$, tel que si $A$ et $B$ sont
les opérations dans $V$ et $W$, celle dans
$V\otimes W$ soit $A\otimes1+1\otimes B$. $\unsure{\mathfrak g|g}$ est une
algèbre de Lie \inserted{topologique complète}, (et $=$ une $p$-algèbre de
Lie si on est en car. $p\ne0$)
\margin{voir plus bas une étude plus \unsure{complète}}
% slanted

% archive sheet number: 169
